\begin{comment}
\documentclass[12pt]{article}
\input{header}
\input{header2}
%\begin{document}
\end{comment}

\begin{proof}[Proof of Proposition \ref{prop:ChoProp3}] 

{% begin temporary command space
\newcommand{\termone}{I(X>t)}
\newcommand{\termtwo}{\omdi(X\leq t) \frac{\Snhat\argt}{\Snhat\argX}}
\newcommand{\termthree} {\Snhat\argt}
\newcommand{\estfunclongtrue}{\termone + \termtwo - \termthree}
\newcommand{\bigE}{E_{\unnormdirac\estfunctruesurvlabel\mid\unnormdirac}}
\newcommand{\bigEone}{E_{\unnormdirac\estfunctruesurvlabel\mid\unnormdirac = 1}}
\newcommand{\smallE}{E_{\unnormdirac}}
By hypothesis, $\norm{\esteqlabelSnhat}_L = \suptzerototauh \bigg| \PdotSnhat \bigg| \to 0$. We want to show that $\PdotSnhat \to 0$ uniformly over $\histlower \in \history$ and $t \in [0, \tau]$ implies $
\sup_{t,\histlower} u_n(t \mid \histlower) \to 0$ for $u_n$ defined below. We know 
\begin{align}
\label{eq:three_components}
\PdotSnhat = \PopP_{\cdot\mid \histlower} \bigg\{ 
\termone
+\termtwo
-\termthree
\bigg\}.
\end{align}
We individually examine each term in \cref{eq:three_components}.
%first term
For the first term, because $T\perp C$, we have 
\begin{align}
\label{eq:component1}
\Pd \left( \termone \right) &= E\left( \termone \mid \heh \right) = \Pr(X>t\mid \heh ) = \Pr(T>t, C>t\mid \heh ) \notag \eqnspace
%component1
&= \Pr(T>t\mid \heh )\Pr(C>t\mid \heh ) = \truesurv(t\mid \heh)\left( 1 - G(t\mid \heh) \right). \tag{\ref{eq:three_components}.1}
\end{align}
} %end tmp command space
{ % temporary command space
\newcommand{\indx}{I(X\leq t)}
\newcommand{\fracX}{\frac{\Snhat\argt}{\Snhat\argX}}
\newcommand{\fracC}{\frac{\Snhat\argt}{\Snhat\argC}}
\newcommand{\fracCheh}{\frac{\Snhat(t\mid \heh)}{\Snhat(C\mid\heh)}}
\newcommand{\argcheh}{(c\mid\heh)}
\newcommand{\argtheh}{(t\mid\heh)}
\newcommand{\arguheh}{(u\mid\heh)}
\newcommand{\intc}{\int_c^\infty}
For the second term,
\begin{align*}
\label{eq:component2}
\Pd \left\{ \omd \indx \fracX \right\} &= E\left( \omd \indx \fracX \bigg| \heh \right) %\notag \eqnspace
\end{align*}
\begin{align}
&= E\left( I(C\leq t) \fracC \bigg| \heh, \delta = 0 \right) \notag \eqnspace
&=\intzerotot \, \intc \fracCheh \, f_{T,C}(u,c \mid \heh )\, du \, dc \notag \eqnspace
&=\intzerotot \frac{1}{\Snhat\argcheh} \intc \Snhat\argtheh f_T\arguheh \, du \, f_C\argcheh \, dc \notag \eqnspace
%component2
&= \intzerotot \frac{f_C\argcheh\truesurv\argcheh}{\Snhat\argcheh}\, dc \, \Snhat\argtheh, \tag{\ref{eq:three_components}.2}
\end{align}
\begin{comment}
\Pd \left\{ \omd &\indx \fracX \right\} = E\left( \omd \indx \fracX \bigg| \heh \right) \notag \eqnspace
&= E\left( I(C\leq t) \fracC \bigg| \heh, \delta = 0 \right) \notag \eqnspace
&=\intzerotot \, \intc \fracCheh \, f_{T,C}(u,c \mid \heh )\, du \, dc \notag \eqnspace
&=\intzerotot \frac{1}{\Snhat\argcheh} \intc \Snhat\argtheh f_T\arguheh \, du \, f_C\argcheh \, dc \notag \eqnspace
%component2
&= \intzerotot \frac{f_C\argcheh\truesurv\argcheh}{\Snhat\argcheh}\, dc \, \Snhat\argtheh, \tag{\ref{eq:three_components}.2}
\end{comment}
where the last equality holds because \begin{align*}
\intc \Snhat\argtheh \, f_T\arguheh \, du = \Snhat\argtheh \, F_T\argu \bigg|^\infty_c &= \Snhat\argtheh \, (1-F_T\argc) \eqnspace &= \Snhat\argtheh \, (1- \truesurv\argc).
\end{align*}
The last term is
\begin{align}
\label{eq:component3}
\Pd \{ -\Snhat\argt\} 
%componenet3
= -\Snhat\argtheh.
\tag{\ref{eq:three_components}.3}
\end{align}
\newcommand{\fracsn}{\frac{\Snhat\argtheh}{\Snhat\argtheh}}
\newcommand{\fracSoSn}{\frac{\truesurv\argtheh}{\Snhat\argtheh}}
Combining \cref{eq:component1}, \cref{eq:component2}, \cref{eq:component3}, we have
\begin{align*}
\PdotSnhat &= 
%component1
\truesurv(t\mid \heh)\left( 1 - G(t\mid \heh) \right) \eqnspace
%component2
&\quad\quad\quad+ \intzerotot \frac{f_C\argcheh\truesurv\argcheh}{\Snhat\argcheh}\, dc \, \Snhat\argtheh
%component3
-\Snhat\argtheh \eqnspace
&= \truesurv(t\mid \heh) -\Snhat\argtheh 
-  \truesurv(t\mid \heh) G(t\mid \heh) \eqnspace
&\quad\quad\quad+ \Snhat\argtheh\intzerotot \frac{f_C\argcheh\truesurv\argcheh}{\Snhat\argcheh}\, dc
 \eqnspace
&= \fracsn \truesurv(t\mid \heh) -\Snhat\argtheh \\
&\quad \quad \quad  
-  \fracsn \truesurv(t\mid \heh) G(t\mid \heh) \\
&\quad \quad \quad + \Snhat\argtheh G\argtheh - \Snhat\argtheh G\argtheh
\eqnspace
&\quad\quad\quad+\Snhat\argtheh\intzerotot \frac{\truesurv\argcheh}{\Snhat\argcheh}\, dG\argcheh \eqnspace
&= \Snhat\argtheh \bigg\{ \fracSoSn - 1 -\left[ \fracSoSn - 1 \right] G\argtheh \\
&\quad \quad \quad \quad \quad \quad \quad  \quad \quad \quad 
- G\argtheh + \intzerotot \frac{\truesurv\argcheh}{\Snhat\argcheh}\, dG\argcheh \bigg\}.
\end{align*}
\newcommand{\en}{\epsilon_n}
\newcommand{\un}{u_n}
Now, letting $\en\argtheh =  \fracSoSn - 1$, and noting that
\begin{align*}
\intzerotot &\frac{\truesurv\argcheh}{\Snhat\argcheh}\, dG\argcheh - G\argtheh \eqnspace &=\intzerotot \frac{\truesurv\argcheh}{\Snhat\argcheh}\, dG\argcheh - \intzerotot 1 \, dG\argcheh \eqnspace
&=\intzerotot \left[ \frac{\truesurv\argcheh}{\Snhat\argcheh}\, - 1 \right] \, dG\argcheh =\intzerotot \en\argcheh \, dG\argcheh,
\end{align*} 
we have $\PdotSnhat$
\begin{align*}
 &= \Snhat\argtheh \bigg\{ \en\argtheh (1 - G\argtheh) \eqnspace
&\quad\quad\quad\quad\quad\quad\quad\quad+ \intzerotot \frac{\truesurv\argcheh}{\Snhat\argcheh}\, dG\argcheh - G\argtheh  \bigg\} \eqnspace
&= \Snhat\argtheh \bigg\{ \en\argtheh (1 - G\argtheh) + \intzerotot \en\argcheh\, dG\argcheh  \bigg\}. %\eqnspace
%&\equiv \Snhat\argtheh \un\argtheh.
\end{align*}
Setting $\un \equiv \en\argtheh (1 - G\argtheh) + \intzerotot \en\argcheh\, dG\argcheh$, and taking the derivative on both sides, we get %using product rule
\begin{equation*}
\begin{aligned}
d\un\argtheh &= d\bigg[ \en\argtheh (1 - G\argtheh) + \intzerotot \en\argcheh\, dG\argcheh \bigg] \eqnspace
&=d\en\argt(1 - G\argtheh) + \bigg(d(1 - G\argtheh)\bigg) \en\argtheh \eqnspace
&\quad \quad + \en\argtheh dG\argtheh \eqnspace
&=d\en\argt(1 - G\argtheh) - dG\argtheh \, \en\argtheh \eqnspace
&\quad \quad + \en\argtheh dG\argtheh \eqnspace
&=d\en\argt(1 - G\argtheh) \eqnspace
\implies d\en\argtheh &= \frac{d\un\argt}{1 - G\argtheh}.
\end{aligned}
\end{equation*}
Using integration by parts, we have
\begin{align*}
\en\argtheh &= \intzerotot \frac{d\un\argt}{1 - G\argtheh} \eqnspace
&= \frac{\un\arguheh}{1-G\arguheh}\bigg|^t_0 - \intzerotot \frac{\un\arguheh \, dG\arguheh}{(1-G\arguheh)^2} \eqnspace
&= \frac{\un\argtheh}{1-G\argtheh} - \intzerotot \frac{\un\arguheh \, dG\arguheh}{(1-G\arguheh)^2}. 
\end{align*}
By Assumption \ifnum\jasa=1 4
 \else \ref{assump:CuiAssump1}\fi, $1-G\argtheh \ge 1- G(\tau\mid\heh) > 0$. Thus, $\un\argtheh \to 0$ implies $\en\argtheh \to 0$. Since $\un\argtheh \to 0$ by hypothesis, we have $\fracSoSn - 1 \to 0$. Thus, $\norm{\Snhat - \truesurv}_\infty \to 0$.
} %end tmp command space

\end{proof}
%\end{document}
